\chapter{Activation}
\label{ch:activation}

\chapterbrief
  {Take a card from ``something is on the antenna'' to ``I know its UID and it is ready for commands''.}
  {Value objects and immutability.}
  {A real UID, SAK and ATS printed from your card.}

\section{The problem: the card's state machine}

A contactless card is a state machine, and you must walk it in order.

\begin{center}
\begin{tikzpicture}[font=\footnotesize\sffamily, node distance=11mm,
   st/.style={draw=brewernavy, rounded corners=2pt, minimum width=26mm,
              minimum height=7.5mm, inner sep=3pt},
   lbl/.style={font=\scriptsize\sffamily, color=breweramber!80!black}]
  \node[st] (off) {POWER-OFF};
  \node[st, below=of off] (idle) {IDLE};
  \node[st, below=of idle] (ready) {READY};
  \node[st, below=of ready] (active) {ACTIVE};
  \node[st, below=of active, draw=brewergreen] (isodep) {ISO-DEP};
  \draw[-{Stealth[length=4pt]}, brewernavy] (off) -- node[right, lbl] {enters the field} (idle);
  \draw[-{Stealth[length=4pt]}, brewernavy] (idle) -- node[right, lbl] {REQA / WUPA} (ready);
  \draw[-{Stealth[length=4pt]}, brewernavy] (ready) -- node[right, lbl] {anticollision + SELECT} (active);
  \draw[-{Stealth[length=4pt]}, brewernavy] (active) -- node[right, lbl] {RATS} (isodep);
  \draw[-{Stealth[length=4pt]}, brewergrey, dashed]
        (isodep.west) -- ++(-16mm,0) |- node[left, lbl, pos=0.25] {leaves the field} (off.west);
  \node[font=\scriptsize\itshape, color=brewergreen, right=3mm of isodep]
       {you can finally send APDUs};
  \node[font=\scriptsize\itshape, color=brewergrey, right=3mm of ready]
       {Chapter~\ref{ch:reader} got you here};
\end{tikzpicture}
\end{center}

\section{Anticollision, and why it exists}
\index{anticollision}

Put two cards on the antenna and both answer at once --- the radio signals
collide and you get mush. ISO-14443 solves this with a bit-by-bit tree walk: the
reader asks ``everyone whose UID starts with these bits, answer'', and narrows
until exactly one card remains.

We implement the \textbf{simple case}: ask for the whole UID at once, and if the
reply's checksum is intact, exactly one card answered. If two cards collide, the
checksum fails, we return empty, and the user is told to present one card. That
is honest engineering for a coffee machine.

\paragraph{BCC.} The block check character is
\code{uid[0] \^{} uid[1] \^{} uid[2] \^{} uid[3]}. XOR the four bytes; the result
must equal the fifth. A cheap integrity check that catches both RF corruption
and undetected collisions.

\subsection{Cascade levels, and the \texttt{0x88} tag}

UIDs come in 4-byte and 7-byte flavours, but an anticollision reply is fixed at
five bytes (four data + one checksum). A 7-byte UID therefore needs two rounds:

\begin{center}
\begin{tikzpicture}[font=\scriptsize\sffamily, x=8mm, y=7mm]
  % level 1
  \node[font=\footnotesize\sffamily, anchor=west, color=brewergrey] at (-0.2,2.35)
       {Level 1 (SEL \hex{0x93}) reply};
  \foreach \i/\lbl/\c in {0/{\hex{88}}/brewerred!12, 1/{u0}/breweramber!14,
                          2/{u1}/breweramber!14, 3/{u2}/breweramber!14,
                          4/{BCC}/black!7}{
     \fill[\c] (\i,1.3) rectangle ++(1,1);
     \draw[brewergrey] (\i,1.3) rectangle ++(1,1);
     \node at (\i+0.5,1.8) {\lbl};
  }
  \node[color=brewerred, font=\scriptsize\itshape, anchor=west] at (5.25,1.8)
       {cascade tag --- NOT part of the UID};
  % level 2
  \node[font=\footnotesize\sffamily, anchor=west, color=brewergrey] at (-0.2,0.85)
       {Level 2 (SEL \hex{0x95}) reply};
  \foreach \i/\lbl in {0/{u3}, 1/{u4}, 2/{u5}, 3/{u6}}{
     \fill[breweramber!14] (\i,-0.2) rectangle ++(1,1);
     \draw[brewergrey] (\i,-0.2) rectangle ++(1,1);
     \node at (\i+0.5,0.3) {\lbl};
  }
  \fill[black!7] (4,-0.2) rectangle ++(1,1);
  \draw[brewergrey] (4,-0.2) rectangle ++(1,1);
  \node at (4.5,0.3) {BCC};
  % result
  \node[font=\footnotesize\sffamily, anchor=west, color=brewergrey] at (-0.2,-0.75)
       {Real UID = 7 bytes};
  \foreach \i/\lbl in {0/{u0}, 1/{u1}, 2/{u2}, 3/{u3}, 4/{u4}, 5/{u5}, 6/{u6}}{
     \fill[brewergreen!12] (\i,-1.9) rectangle ++(1,1);
     \draw[brewergreen!60] (\i,-1.9) rectangle ++(1,1);
     \node at (\i+0.5,-1.4) {\lbl};
  }
\end{tikzpicture}
\end{center}

The SAK from level 1 comes back with \textbf{bit 2 (\hex{0x04}) set},
confirming ``UID incomplete, do the next level''. DESFire EV2/EV3 always has a
7-byte UID, so you \emph{will} exercise this path.

\begin{trap}{Forgetting to skip the \hex{0x88}}
Keep the cascade tag and you get a 7-byte UID that is wrong in its first byte
--- and it looks completely plausible. You will only notice when it does not
match what your phone reports, which might be days later.
\end{trap}

\paragraph{RATS.} Request Answer To Select (\hex{0xE0} \hex{0x50}) is the final
handshake: ``let us speak ISO-14443-4''. The card replies with its \textbf{ATS},
describing its capabilities and maximum frame size. The \hex{0x50} parameter is
FSDI${}=5$ (we accept up to 64-byte frames) and CID${}=0$.

\section{Build it}

\subsection{\texttt{ActivatedCard} --- a value object}
\index{value object}\index{immutability}

\begin{jcode}
package com.midnightbrewer.hardware;

/**
 * What activation produced: a UID, a SAK and an ATS.
 *
 * A VALUE OBJECT -- it holds data, has no behaviour that changes it, and is
 * IMMUTABLE: every field is final and there is no setter. Once built, it can
 * never disagree with itself.
 *
 * Why bother, instead of returning byte[][] or three out-parameters?
 *   - the three fields arrive together and belong together
 *   - the compiler enforces that you cannot forget one
 *   - it has a NAME, so signatures read as intent:
 *       ActivatedCard activate()      vs      byte[][] activate()
 */
public final class ActivatedCard {

    private final byte[] uid;
    private final int sak;
    private final byte[] ats;

    public ActivatedCard(byte[] uid, int sak, byte[] ats) {
        // DEFENSIVE COPY on the way in: if we stored the caller's array
        // directly, they could mutate our state afterwards and we would
        // silently change. Copy, and we own it.
        this.uid = uid.clone();
        this.sak = sak;
        this.ats = ats.clone();
    }

    /** Defensive copy on the way OUT too, for the same reason. */
    public byte[] uid() { return uid.clone(); }

    public int sak() { return sak; }

    public byte[] ats() { return ats.clone(); }

    /** The UID as spaced uppercase hex, e.g. "04 76 18 62 7B 20 90". */
    public String uidHex() {
        // TODO: StringBuilder, String.format("%02X ", b) per byte, trim().
        return null;
    }
}
\end{jcode}

\begin{concept}{Three decisions, and why each one}
\textbf{\code{final class}} --- nobody may subclass it. For a value object that
is right: there is no such thing as a ``special kind of ActivatedCard'', and
allowing subclasses would let someone override \code{uid()} to return something
else, breaking the guarantee.

\textbf{Defensive copies both ways} --- the part people skip. Arrays in Java are
references; handing out \code{this.uid} hands out a pointer to your internals,
and the caller can scribble on it. \code{clone()} costs seven bytes and buys you
a class that genuinely cannot be corrupted from outside.

\textbf{Immutability} --- an \code{ActivatedCard} is \emph{a fact about a
moment}: ``at 14:32 this card, with this UID, was activated''. Facts do not
change. Making it immutable means you can pass it anywhere, store it, log it,
hand it to another thread, and never wonder whether something modified it.
\end{concept}

\subsection{The activation sequence}

Add these constants to \code{Iso14443Reader}:

\begin{jcode}
    /** Anticollision/SELECT command bytes for the two cascade levels. */
    private static final byte SEL_CL1 = (byte) 0x93;
    private static final byte SEL_CL2 = (byte) 0x95;

    /** NVB byte that turns an anticollision into a SELECT (full UID follows). */
    private static final byte SELECT_NVB = 0x70;

    /** SAK bit 3: "UID incomplete, do the next cascade level". */
    private static final int CASCADE_BIT = 0x04;

    /** RATS: request ISO-14443-4. PARAM 0x50 = FSDI 5, CID 0. */
    private static final byte RATS_CMD   = (byte) 0xE0;
    private static final byte RATS_PARAM = 0x50;

    /** The tag prefixing a 7-byte UID's first anticollision reply. */
    private static final byte CASCADE_TAG = (byte) 0x88;

    /** Anticollision reply is 4 UID/CT bytes + 1 BCC. */
    private static final int ANTICOLL_LENGTH = 5;
\end{jcode}

And the methods:

\begin{jcode}
    // =================================================================
    // anticollision -- ask the card for (part of) its UID.
    //
    //   C: serNum[0] = 0x93; serNum[1] = 0x20;   // 0x20 NVB = "send it all"
    //      ToCard(PCD_TRANSCEIVE, serNum, 2, ...);
    //      // verify serNum[0]^[1]^[2]^[3] == serNum[4]
    //
    // Takes selCode so it works at BOTH levels. Full bytes, so
    // bitsInLastByte = 0. Returns the 5-byte reply, or empty if invalid.
    // =================================================================
    public byte[] anticollision(byte selCode) throws SpiException {
        // TODO 1: transceive { selCode, 0x20 } with 0 bits-in-last-byte.
        // TODO 2: a valid reply is exactly ANTICOLL_LENGTH bytes; else empty.
        // TODO 3: verify the BCC (XOR of the first four == the fifth);
        //         if it fails the read was garbled -> return empty.
        // TODO 4: return the reply.
        return new byte[0];
    }

    // =================================================================
    // select -- send the UID back, receive the SAK.
    //
    //   C: buffer[0]=selCode; buffer[1]=0x70;      // 0x70 NVB = SELECT
    //      memcpy(&buffer[2], anticoll, 5);        // 4 UID + BCC
    //      CalculateCRC(buffer, 7, &buffer[7]);    // -> 9 bytes
    //      ToCard(...);   // reply = SAK(1) + CRC(2)
    //
    // @return the SAK byte, or -1 if the reply was malformed.
    // =================================================================
    public int select(byte selCode, byte[] anticoll) throws SpiException {
        // TODO 1: build a 7-byte frame { selCode, SELECT_NVB, ...anticoll(5) }.
        // TODO 2: byte[] crc = calculateCrc(frame7);
        //         build a 9-byte frame: the 7, then crc[0], crc[1].
        // TODO 3: byte[] reply = transceive(frame9, 0);
        // TODO 4: if reply.length < 1 return -1, else return reply[0] & 0xFF.
        return -1;
    }

    // =================================================================
    // rats -- ask to speak ISO-14443-4; the card returns its ATS.
    //   C: req[0]=0xE0; req[1]=0x50; CalculateCRC(req,2,&req[2]); ToCard(...)
    // =================================================================
    public byte[] rats() throws SpiException {
        // TODO: header { RATS_CMD, RATS_PARAM }, append its CRC, transceive.
        return new byte[0];
    }

    // =================================================================
    // activate() -- the whole sequence, as another TEMPLATE METHOD.
    //
    // 'final' so a subclass cannot reorder the protocol.
    //   level 1: anticollision -> select -> SAK
    //   if the SAK's cascade bit is set, the UID is 7 bytes: do level 2
    //   then RATS
    // and assembles the real UID:
    //   - 4-byte card: ac1[0..3]
    //   - 7-byte card: ac1[1..3]  (SKIP the 0x88 tag)  +  ac2[0..3]
    // =================================================================
    public final ActivatedCard activate() throws SpiException {
        // TODO 1: ac1 = anticollision(SEL_CL1); if empty -> throw SpiException.
        //         sak1 = select(SEL_CL1, ac1);
        // TODO 2: if ((sak1 & CASCADE_BIT) == 0) -> 4-byte UID:
        //         uid = ac1[0..3], return new ActivatedCard(uid, sak1, rats());
        // TODO 3: else 7-byte UID: ac2 = anticollision(SEL_CL2);
        //         sak2 = select(SEL_CL2, ac2);
        //         uid = { ac1[1], ac1[2], ac1[3], ac2[0..3] }
        //         return new ActivatedCard(uid, sak2, rats());
        return null;
    }
\end{jcode}

\section{Verify on hardware}

\begin{jcode}
if (reader.isCardPresent()) {              // wake the card into READY first
    ActivatedCard card = reader.activate();
    System.out.println("UID  = " + card.uidHex());
    System.out.printf("SAK  = 0x%02X%n", card.sak());
    System.out.println("ATS  = " + hex(card.ats()));
}
\end{jcode}

\begin{verifybox}
For a DESFire EV3:
\begin{trace}
UID  = 04 76 18 62 7B 20 90
SAK  = 0x20
ATS  = 06 75 77 81 02 80 02 F0
\end{trace}
\end{verifybox}

\paragraph{How to read that output.}

\begin{itemize}
  \item \textbf{UID starts \hex{04}} --- NXP's manufacturer byte. Seven bytes
        total, so the cascade path ran and you skipped the tag correctly.
  \item \textbf{\hex{SAK = 0x20}} --- bit 5 set means ``ISO-14443-4
        compliant''. This is the card telling you \emph{I speak APDUs}. A MIFARE
        Classic answers \hex{0x08}; if you see that, you have the wrong card and
        nothing after this chapter will work.
  \item \textbf{ATS} --- \hex{06} is its length, and the \hex{75} encodes
        FSCI${}=5$, i.e.\ a 64-byte maximum frame.
\end{itemize}

\begin{trap}{The FSCI nibble}
FSCI lives in one nibble of the ATS, and it is genuinely easy to read the wrong
one --- the original C in this project had exactly that bug.

Read the high nibble instead of the low and you conclude FSCI${}=7$, i.e.\
``128-byte frames''. You then send a 100-byte APDU, and the card silently
ignores it. Symptom: everything works until one particular command, which times
out forever.
\end{trap}

\section{Traps}

\begin{traptable}
UID is 7 bytes but wrong in byte 0 & Did not skip the \hex{0x88} cascade tag. \\
\code{select()} returns $-1$ & CRC wrong. Check \reg{ModeReg}$=$\hex{0x3D} (preset \hex{0x6363}). \\
Anticollision always returns empty & Card not in READY --- call \code{isCardPresent()} first. \\
\hex{SAK = 0x08} & It is a MIFARE Classic, not a DESFire. \\
RATS returns empty & Card does not support 14443-4, or you sent RATS twice. \\
Works, then fails on retry & Card left in ACTIVE. Power-cycle the field per transaction. \\
\end{traptable}

\section{Commit}

\begin{shell}
git add -A
git commit -m "M5: activation -- anticollision, SELECT, cascade, RATS, ActivatedCard"
\end{shell}
